\subsection{微分}

\subsubsection{微分係数と導関数}

関数 $y=f(x)$ において、$x$ を $\Delta x$ だけ変化させたときの
$y$ の変化量は $\Delta y=f(x+\Delta x)-f(x)$ である。その比
\[
 \frac{\Delta y}{\Delta x}
 =\frac{f(x+\Delta x)-f(x)}{\Delta x}
\]
を平均変化率という。$\Delta x\to0$ の極限が存在するとき、
\begin{equation}
 \boxed{f'(x)=\lim_{\Delta x\to0}
 \frac{f(x+\Delta x)-f(x)}{\Delta x}}
\end{equation}
を $f$ の\textbf{導関数}という。$x=a$ における値 $f'(a)$ は
微分係数であり、曲線 $y=f(x)$ の点 $(a,f(a))$ における接線の傾きである。

\begin{center}
\begin{tikzpicture}[>=Stealth,scale=1.15]
 \draw[->] (-.4,0)--(5.2,0) node[right] {$x$};
 \draw[->] (0,-.3)--(0,3.5) node[above] {$y$};
 \draw[thick,domain=-.1:4.5,samples=80,smooth]
   plot(\x,{.16*(\x-1.2)^2+.5});
 \coordinate (P) at (2,{.16*(2-1.2)^2+.5});
 \coordinate (Q) at (4,{.16*(4-1.2)^2+.5});
 \fill (P) circle (1.5pt) node[below left] {$P$};
 \fill (Q) circle (1.5pt) node[above left] {$Q$};
 \draw[dashed] (P)--(2,0) node[below] {$x$};
 \draw[dashed] (Q)--(4,0) node[below] {$x+\Delta x$};
 \draw[blue,thick] ($(P)!-.35!(Q)$)--($(P)!1.25!(Q)$)
   node[right] {割線};
 \draw[red,thick] (.5,.22)--(3.8,1.07) node[right] {接線};
\end{tikzpicture}

\small 図\thefigure\quad 割線は $\Delta x\to0$ で接線に近づく
\end{center}

高階導関数は
\[
 f''(x)=\frac{d^2f}{dx^2},\qquad
 f^{(n)}(x)=\frac{d^nf}{dx^n}
\]
と表す。また $x=a$ における値は
$f^{(n)}(a)=\left.\dfrac{d^nf}{dx^n}\right|_{x=a}$ と書く。

\subsubsection{定義から導く基本公式}

\paragraph{べき関数}
二項定理を用いると
\begin{align*}
 \frac{(x+\Delta x)^n-x^n}{\Delta x}
 &=nx^{n-1}+\binom n2x^{n-2}\Delta x+\cdots+(\Delta x)^{n-1}.
\end{align*}
したがって
\begin{equation}
 \boxed{\frac{d}{dx}x^n=nx^{n-1}}
\end{equation}
である。

\paragraph{指数関数}
自然対数の底 $e$ は
\[
 \lim_{h\to0}\frac{e^h-1}{h}=1
\]
を満たすので、
\[
 \frac{d}{dx}e^x
 =e^x\lim_{h\to0}\frac{e^h-1}{h}=e^x.
\]

\paragraph{三角関数}
加法定理から
\begin{align*}
 \frac{\sin(x+h)-\sin x}{h}
 &=\sin x\,\frac{\cos h-1}{h}
   +\cos x\,\frac{\sin h}{h}.
\end{align*}
ここで
\[
 \lim_{h\to0}\frac{\sin h}{h}=1,\qquad
 \lim_{h\to0}\frac{\cos h-1}{h}=0
\]
であるから、$(\sin x)'=\cos x$ を得る。

最初の極限は単位円上の面積比較から示せる。$0<h<\pi/2$ について
\[
 \frac12\sin h<\frac12h<\frac12\tan h
\]
なので、$\cos h<\sin h/h<1$ である。$h\to0$ として挟み撃ちの
原理を用いればよい。また、余弦の加法定理を用いると
\begin{align*}
 \frac{\cos(x+h)-\cos x}{h}
 &=\cos x\,\frac{\cos h-1}{h}
   -\sin x\,\frac{\sin h}{h}\\
 &\longrightarrow-\sin x
\end{align*}
なので、$(\cos x)'=-\sin x$ である。

\begin{center}
\begin{tikzpicture}[>=Stealth,scale=2]
 \coordinate (O) at (0,0);
 \coordinate (A) at (1,0);
 \coordinate (B) at ({cos(35)},{sin(35)});
 \coordinate (T) at (1,{tan(35)});
 \draw (O)--(A) arc[start angle=0,end angle=35,radius=1]--cycle;
 \draw (A)--(T)--(O);
 \draw[dashed] (B)--({cos(35)},0);
 \draw[thick,blue] (O)--(B);
 \draw[thick,red] (O)--(T);
 \draw[->] (.35,0) arc[start angle=0,end angle=35,radius=.35]
   node[midway,right] {$h$};
 \node[below] at (A) {$A$};
 \node[above left] at (B) {$P$};
 \node[right] at (T) {$T$};
 \node[left] at (O) {$O$};
 \node[below] at (.5,0) {$1$};
 \node[right] at ({cos(35)},{sin(35)/2}) {$\sin h$};
 \node[right] at (1,{tan(35)/2}) {$\tan h$};
\end{tikzpicture}

\small 単位円上の三角形と扇形の面積比較
\end{center}

\paragraph{対数関数}
$x>0$ とし、$h/x=1/n$ とおくと
\begin{align*}
 \frac{\log(x+h)-\log x}{h}
 &=\frac1x\frac{x}{h}\log\left(1+\frac hx\right)\\
 &=\frac1x\log\left(1+\frac1n\right)^n\longrightarrow\frac1x.
\end{align*}
したがって $(\log x)'=1/x$ である。ここでは $\log$ は自然対数を表す。

\subsubsection{微分法則}

微分可能な関数 $f,g$ について
\begin{align}
 (f+g)'&=f'+g',\\
 (fg)'&=f'g+fg',\\
 \left(\frac fg\right)'&=\frac{f'g-fg'}{g^2}\qquad(g\ne0)
\end{align}
が成り立つ。積の微分は、途中に $f(x)g(x+h)$ を挿入して
\begin{align*}
 \frac{f(x+h)g(x+h)-f(x)g(x)}h
 &=\frac{f(x+h)-f(x)}h\,g(x+h)\\
 &\quad+f(x)\frac{g(x+h)-g(x)}h
\end{align*}
と分け、$h\to0$ とすれば得られる。商の微分は
$f/g=f\cdot g^{-1}$ とみなし、
\[
 (g^{-1})'=-\frac{g'}{g^2}
\]
を積の微分と組み合わせれば導ける。

合成関数 $f(g(x))$ について、$\Delta g=g(x+\Delta x)-g(x)$ とおくと
\[
 \frac{f(g+\Delta g)-f(g)}{\Delta x}
 =\frac{f(g+\Delta g)-f(g)}{\Delta g}\frac{\Delta g}{\Delta x}.
\]
極限を取ることにより、\textbf{連鎖律}
\begin{equation}
 \boxed{\frac{d}{dx}f(g(x))=f'(g(x))g'(x)}
\end{equation}
を得る。たとえば
\[
 \frac{d}{dx}\sqrt{1+x^2}=\frac{x}{\sqrt{1+x^2}},\qquad
 \frac{d}{dx}\tan x=\frac1{\cos^2x}.
\]

\subsubsection{停留点・極値・変曲点}

$f'(a)=0$ となる点 $x=a$ を\textbf{停留点}という。停留点の前後で
$f'$ が正から負へ変われば極大、負から正へ変われば極小である。
$f''(a)<0$ なら極大、$f''(a)>0$ なら極小と判定できるが、
$f''(a)=0$ の場合は高階導関数や $f'$ の符号を調べる必要がある。
また、曲線の凹凸が変わる点を\textbf{変曲点}という。$f''(a)=0$ は
候補を与えるだけで、前後で $f''$ の符号が変わることを確認する。

\begin{center}
\begin{tikzpicture}[>=Stealth,scale=.9]
 \draw[->] (-4.6,0)--(4.8,0) node[right] {$x$};
 \draw[->] (0,-2.2)--(0,2.8) node[above] {$y$};
 \draw[thick,smooth] plot coordinates {
 (-4,-1)(-3.2,1.4)(-2.5,2)(-1.8,1.2)(-.8,.3)
 (0,-.2)(.8,-.25)(1.8,-.8)(2.7,-1.6)(3.5,-.8)(4.3,1.3)};
 \foreach \x/\y/\lab in {-2.5/2/A,.8/-.25/C,2.7/-1.6/E}
   {\fill (\x,\y) circle (1.7pt); \draw[dashed] (\x,\y)--(\x,0)
    node[above right] {$x_{\lab}$};}
 \foreach \x/\y/\lab in {-1.55/.92/B,1.75/-.76/D}
   {\fill (\x,\y) circle (1.4pt); \draw[densely dashed,gray] (\x,\y)--(\x,0)
    node[below] {$x_{\lab}$};}
\end{tikzpicture}

\small 極大点 $x_A$、極小点 $x_E$、停留点 $x_C$ と
変曲点 $x_B,x_D$
\end{center}

\paragraph{演習}
\[
 f(x)=\frac15x^5-\frac34x^4+\frac13x^3+\frac32x^2-2x
\]
の増減、極値および変曲点を調べよ。

\paragraph{解答}
\begin{align*}
 f'(x)
 &=x^4-3x^3+x^2+3x-2\\
 &=(x+1)(x-1)^2(x-2)
\end{align*}
である。よって
\[
\begin{array}{c|ccccccc}
x&(-\infty,-1)&-1&(-1,1)&1&(1,2)&2&(2,\infty)\\ \hline
f'(x)&+&0&-&0&-&0&+
\end{array}
\]
となる。したがって $x=-1$ で極大、$x=2$ で極小となり、
$x=1$ は停留点だが極値ではない。関数値は
\[
 f(-1)=\frac{133}{60},\qquad
 f(1)=-\frac{43}{60},\qquad
 f(2)=-\frac{14}{15}
\]
である。

また
\begin{align*}
 f''(x)
 &=4x^3-9x^2+2x+3\\
 &=(x-1)(4x^2-5x-3)\\
 &=4(x-1)
 \left(x-\frac{5-\sqrt{73}}8\right)
 \left(x-\frac{5+\sqrt{73}}8\right).
\end{align*}
各零点の前後で $f''$ の符号が変わるので、変曲点の $x$ 座標は
\[
 \boxed{x=\frac{5-\sqrt{73}}8,\quad1,\quad
        \frac{5+\sqrt{73}}8}
\]
である。
